
\section{多值截面与可数链集的选择}
\label{sec:7-2}

现在可以利用第~7.1 节构造的仓西结构，
建立基本链，
从而严格说明前几章所用的各项构造。
先概述本节安排。
第~7.2.1--7.2.3 小节
说明如何构造满足适当相容条件的多值截面；
本节余下部分将详细实施这一构造。
除
\textbf{[FuOn99II]}
第 1 章确立的仓西结构方法外，
这里还有两个新问题，
分别在第~7.2.2、7.2.3 小节讨论。
这两点颇为微妙，
相关文献中曾出现各种错误、
缺口与混淆；
我们将在本节作详细解释。

\subsection{对角线处的横截性}
\label{subsec:7-2-1}
\liftIndex{I237}%
首先指出，
除非扰动
仓西映射
$s_{k,\beta_0}$（回忆 $\beta_0=0$），
下述断言并不成立：
\begin{quote}
\itshape
“在空间
$\mathcal M^{\mathrm{main}}_{k+1}(\beta)$
的仓西结构上，
存在一族多值截面
$s_{k+1,\beta}$，
与命题~\ref{prop:7-1-2}
中的同构相容。”
\end{quote}

说明尝试证明此断言时会遇到的困难。
设
$\beta_1\in\pi_2(M,L)$
满足
$\beta_1\ne0$、
$\mathcal M^{\mathrm{main}}_1(\beta_1)
\ne\varnothing$，
且对任意满足
$0<\omega(\beta)<\omega(\beta_1)$
的 $\beta$，
都有
$\mathcal M^{\mathrm{main}}_1(\beta)
=\varnothing$。
首先选取
$s_{k+1,\beta_1}$。
为简便，
再假设当
$\omega(\beta_1)
\le\omega(\beta)
<2\omega(\beta_1)$
且 $\beta\ne\beta_1$ 时，
$\mathcal M^{\mathrm{main}}_1(\beta)
=\varnothing$。
下一步便是构造
$s_{k+1,2\beta_1}$。
注意
\begin{equation}
\begin{aligned}
S_2\mathcal M^{\mathrm{main}}_1(2\beta_1)
={}&
\mathcal M^{\mathrm{main}}_1(\beta_1)
{}_{\operatorname{ev}_0}
\times_{\operatorname{ev}_1}
\mathcal M^{\mathrm{main}}_3(\beta_0)
{}_{\operatorname{ev}_2}\\
&\quad\times_{\operatorname{ev}_0}
\mathcal M^{\mathrm{main}}_1(\beta_1),
\end{aligned}
\tag{7.2.1}\label{eq:7-2-1}
\end{equation}
其中 $\beta_0=0$。
构造这样的
$s_{k+1,2\beta_1}$
之所以困难，
是因为
式~\eqref{eq:7-2-1}
右端的纤维积\emph{永不横截}，
即使
$\mathcal M^{\mathrm{main}}_1(\beta_1)$
本身已经横截。
原因是链
\[
\operatorname{ev}
\bigl(
\mathcal M^{\mathrm{main}}_1(\beta_1)
\bigr)
\]
除非具有 $L$ 的满维数，
否则不会与自身横截。
单靠在
$\mathcal M^{\mathrm{main}}_1(\beta_1)$
上选取任何刚才所说的扰动，
也不能使它横截。
这里的问题在于，
为了与
命题~\ref{prop:7-1-2}
的同构相容，
式~\eqref{eq:7-2-1}
右端两个
$\mathcal M^{\mathrm{main}}_1(\beta_1)$
副本上必须使用\emph{同一个}扰动。

这个问题已在
定理~\ref{thm:3-4-8}
与命题~\ref{prop:3-4-5}
的证明中出现：
那里形变相交配对，
使它在链层面有定义，
并据此定义 $A_\infty$ 代数。
通常的链层面相交配对
只对奇异链的某个稠密子集
中的元素有意义；
尤其一条链与自身的
相交配对在链层面并无良定义。
代数拓扑学家似乎早已熟知此事。

不过，
若在 $A_\infty$ 代数范畴，
而非微分分次代数范畴中工作，
就可以把奇异链复形
（的一个可数生成子复形）
扰动到另一个
乘法处处有定义的复形。
自 20 世纪六七十年代
有理同伦论发展以来，
代数拓扑学家似乎也已知晓这一点。
这里仍将给出证明，
见定理~\ref{thm:3-4-8}。

先简述思路。
设 $C(L)$ 是
$L$ 上光滑奇异链复形 $S(L)$
的一个可数生成子复形。
选取同恒等映射同痕的
两个微分同胚
$\varphi_0,\varphi_1\colon L\to L$。
再取 $L$ 上光滑奇异单形的
可数集 $\mathcal X_0$。
容易看出，
若 $\varphi_i$ 选得一般，
则对任意
$P,Q\in\mathcal X_0$，
$\varphi_0(P)$
与 $\varphi_1(Q)$ 横截；
这里允许 $P=Q$。
令
\[
\mathfrak m_2(P,Q)
=\varphi_0(P)\cap\varphi_1(Q).
\]
更准确地说，
将此交集作单纯剖分，
并把所得
$\mathfrak m_2(P,Q)$
视作奇异链。
于是对
$\mathcal X_0$
中任意一对元素，
$\mathfrak m_2$
都有定义；
但这一运算不满足结合律，
因为
\[
\begin{aligned}
\mathfrak m_2(\mathfrak m_2(P,Q),R)
&=
\varphi_0\bigl(
\varphi_0(P)\cap\varphi_1(Q)
\bigr)
\cap\varphi_1(R)\\
&\ne
\varphi_0(P)\cap
\varphi_1\bigl(
\varphi_0(Q)\cap\varphi_1(R)
\bigr)\\
&=
\mathfrak m_2(P,\mathfrak m_2(Q,R)).
\end{aligned}
\]
不过，
从
$(P,Q,R)\mapsto
\mathfrak m_2(\mathfrak m_2(P,Q),R)$
到
$(P,Q,R)\mapsto
\mathfrak m_2(P,\mathfrak m_2(Q,R))$
存在链同伦。
把 $\mathfrak m_3$
定义为这两个运算之间的链同伦，
再继续此过程，
便可得到一个
$A_\infty$ 代数。

接着还有另一困难：
需要选取光滑奇异单形的可数集
$\mathcal X(L)$\liftIndex{I238}，
使对任何
$P_1,\ldots,P_k\in\mathcal X(L)$，
$\mathfrak m_k(P_1,\ldots,P_k)$
仍可由
$\mathcal X(L)$ 的元素生成。
困难具体如下。
先固定 $L^3$
中对角线 $\Delta_3$
的扰动，
以定义 $\mathfrak m_2$。
再一般地选取
$\mathcal X_1(L)$，
使任意
$P_1,P_2\in\mathcal X_1(L)$
都满足
$L\times P_1\times P_2$
与 $\Delta_3$ 横截。
于是
$\mathfrak m_2(P_1,P_2)$
有定义。
随后必须把
$\mathfrak m_2(P_1,P_2)$
的三角剖分中出现的
全部光滑奇异单形
加入 $\mathcal X_1(L)$，
得到 $\mathcal X_2(L)$。
这时又须证明，
对
$P_1,P_2\in\mathcal X_2(L)$，
$L\times P_1\times P_2$
仍与 $\Delta_3$ 横截。
若重新一般地选取
$\Delta_3$，
当然可以做到；
但这里
$\Delta_3$ 已先选定，
而 $\mathcal X_2(L)$
恰依赖于它的选择。

因此我们采取如下方法：
让对角线扰动
（以及模空间
$\mathcal M^{\mathrm{main}}_{k+1}(\beta)$
上的多值截面）
\emph{依赖于}
参与纤维积的奇异单形
$P_i$ 的选择。
所以，
如前几章已说明的，
我们直接在纤维积
\[
\mathcal M^{\mathrm{main}}_{k+1}(\beta)
\times\prod_i P_i
\]
上选取多值截面；
\emph{并不是}先在
$\mathcal M^{\mathrm{main}}_{k+1}(\beta)$
上选取多值截面 $s$，
再取纤维积
\[
\mathcal M^{\mathrm{main}}_{k+1}(\beta)^s
\times\prod_i P_i.
\]
本书 2000 年 12 月的预印本
\textbf{[FOOO00]}
曾以另一种方式表述这一点；
这里予以更正。
